\documentclass[10pt,a4paper]{amsart}
\usepackage[margin=19mm]{geometry}
\usepackage{amsmath,amssymb,mathtools}
\usepackage[scr=rsfs,cal=euler]{mathalfa}
\usepackage{xfrac}
\usepackage[unicode,hidelinks,bookmarks=false]{hyperref}
\newcommand{\ie}{i.e.\ }
\newcommand{\cf}{cf.\ }
\newcommand{\resp}{respectively\ }
\newcommand{\Subsection}{\subsection}
\providecommand{\nobreakdash}{\nobreak\hbox{-}}
\setlength{\emergencystretch}{2em}
\multlinegap0pt
\newcommand{\skpt}{\par\smallskip}

\newcommand{\Changel}{\mathcode`l="7160}
\newcommand{\Changelback}{\mathcode`l="716C}
\newcommand{\NoChangel}[1]{%
\expandafter\let\csname old\string#1\endcsname=#1
\let#1=\relax
\newcommand{#1}{\mathop{\mathcode`l="716C\csname old\string#1\endcsname\mathcode`l="7160}\displaylimits}%
}

\NoChangel{\log}
\NoChangel{\ln}
\NoChangel{\lim}
\NoChangel{\limsup}
\NoChangel{\liminf}
\NoChangel{\varlimsup}
\NoChangel{\varliminf}
\NoChangel{\varinjlim}
\NoChangel{\varprojlim}

\newcommand{\sbullet}{{\scriptscriptstyle\bullet}}
\newcommand{\csbullet}{{\raisebox{1pt}{$\sbullet$}}}
\def\mto{\mathchoice{\longmapsto}{\mapsto}{\mapsto}{\mapsto}}
\newcommand{\Psfrac}[2]{(\sfrac{#1}{#2})}
\newcommand{\psfrac}[2]{\sfrac{(#1)}{#2}}
\newcommand{\spfrac}[2]{\sfrac{#1}{(#2)}}
\newcommand{\pspfrac}[2]{\sfrac{(#1)}{(#2)}}
\let\oldbigsqcup\bigsqcup
\renewcommand{\bigsqcup}{\mathop{\textstyle\oldbigsqcup}\displaylimits}

\AtBeginDocument{%
\renewcommand{\labelitemi}{$\sbullet$}
}

\let\faktor\sfrac

\usepackage{tikz,yfonts}
\DeclareFontFamily{U}{mathx}{}
\DeclareFontShape{U}{mathx}{m}{n}{<-> mathx10}{}
\DeclareSymbolFont{mathx}{U}{mathx}{m}{n}
\DeclareMathAccent{\widehat}{0}{mathx}{"70}
\DeclareMathAccent{\widecheck}{0}{mathx}{"71}

\usepackage{tikz-cd}
\usetikzlibrary{decorations.pathmorphing}

\newtheorem{thmx}{Theorem}
\renewcommand{\thethmx}{\Alph{thmx}}
\newtheorem{lemma}{Lemma}[section]
\newtheorem{theorem}[lemma]{Theorem}
\newtheorem{corollary}[lemma]{Corollary}
\newtheorem{proposition}[lemma]{Proposition}
\newtheorem{conjecture}[lemma]{Conjecture}
\newtheorem{question}[lemma]{Question}
\theoremstyle{definition}
\newtheorem{definition}[lemma]{Definition}
\theoremstyle{remark}
\newtheorem{remark}[lemma]{Remark}

\newcommand{\C}{\mathbb{C}}
\newcommand{\D}{\mathbb{D}}
\newcommand{\R}{\mathbb{R}}
\newcommand{\Z}{\mathbb{Z}}
\renewcommand{\P}{{\mathbb P}}

\newcommand{\pslc}{{\mathrm{PSL}_2 (\mathbb{C})}}
\newcommand{\pslr}{{\mathrm{PSL}_2 (\mathbb{R})}}

\newcommand{\cC}{\mathcal{C}}
\newcommand{\cM}{\mathcal{M}}
\newcommand{\cS}{\mathcal{S}}
\newcommand{\cT}{\mathcal{T}}
\newcommand{\cK}{\mathcal{K}}

\DeclareMathOperator{\Int}{int}
\DeclareMathOperator{\Cl}{cl}

\let\phi\varphi
\newcommand{\mate}{\bot \!\! \! {\bot}}

\makeatletter
\DeclareFontFamily{U}{tipa}{}
\DeclareFontShape{U}{tipa}{m}{n}{<->tipa10}{}
\newcommand{\arc@char}{{\usefont{U}{tipa}{m}{n}\symbol{62}}}

\newcommand{\arc}[1]{\mathpalette\arc@arc{#1}}

\newcommand{\arc@arc}[2]{
\sbox0{$\m@th#1#2$}
\vbox{
\hbox{\resizebox{\wd0}{\height}{\arc@char}}
\nointerlineskip
\box0
}
}
\makeatother

\title[Conformal factor-map mating]{Matings, holomorphic correspondences, and a Bers slice: original Theorem3.2}
\author{Mahan Mj and Sabyasachi Mukherjee}
\date{Source: J.\'E.P.--M.12(2025),1445--1502; DOI10.5802/jep.315}
\begin{document}\maketitle
\noindent\textit{Editorial scope.} The counted unit is original Theorem3.2 with all standing orbifold, degree, marked-group and hyperbolic-polynomial hypotheses. Complete original factor-map, normalized conjugacy, Markov/David-extension, invariant Beltrami, holomorphic realization and uniqueness arguments are retained. Finite illustrative pictures and later uniformization/correspondence applications are omitted; all geometric data used here are defined in the retained original text. Author mathematics remains unchanged.
\section{Original orbifold class}\label{intro_sec}
\subsection*{The class of orbifolds}
Before we state the main theorems of the paper, let us describe the general class of orbifolds (equivalently, Fuchsian groups) that are the principal players in the game.
The family of correspondences most extensively studied by Bullett and his collaborators exhibit matings of the modular group $\mathrm{PSL}_2(\Z)$ and quadratic polynomial/rational maps. On the other hand, the conformal matings framework of \cite{MM1} applies to Bowen-Series maps of Fuchsian punctured sphere groups, possibly with an order two elliptic element. In~this paper, we~work with the following collection
of finite volume hyperbolic orbifolds that includes both these as special cases:

$\cS:=$ hyperbolic orbifolds of genus zero with
\begin{enumerate}
\item at least one puncture,
\item at most one order two orbifold point,
\item at most one order $\nu\geq 3$ orbifold point.
\end{enumerate}


\section{Factor Bowen-Series maps}\label{factor_bs_sec}

We will now study Bowen-Series maps associated with appropriate cyclic covers of genus zero orbifolds. These maps have mild discontinuities. However, one can pass to factors of these Bowen-Series maps such that the factors are continuous. The construction of factor Bowen-Series maps is the first key step in the proofs of our main theorems.

\Subsection{Factor Bowen-Series map for a base group}\label{factor_bs_base_subsec}

Let $n, p$ be two positive integers with $np\geq 3$. For $r\in\{1,\dots,n\}$, denote the counter-clockwise arc $\arc{e^{\sfrac{2i\pi(r-1)}{n}}, e^{\sfrac{2i\pi r}{n}}}\subset\mathbb{S}^1$ by $J_r$. Note that $J_1$ is the counter-clockwise arc of $\mathbb{S}^1$ connecting $1$ to $e^{\sfrac{2i\pi}{n}}$, and the various $J_r$ are obtained by rotat\-ing~$J_1$ successively by angle $\sfrac{2\pi}{n}$ about the origin. We~set $\omega:=e^{\sfrac{2i\pi}{n}}$, and \hbox{$M_\omega:\overline{\D}\to\overline{\D}$}, $z\mto \omega z$.

Further, for $r\in\{1,\dots,n\}$, consider the chain of $p$ bi-infinite hyperbolic geodesics
\[
C_{r,s}:=\overline{e^{\Psfrac{2i\pi(r-1)}{n}+\Psfrac{2i\pi(s-1)}{np}}, e^{\Psfrac{2i\pi (r-1)}{n}+\Psfrac{2i\pi s}{np}}},\quad s\in\{1,\dots,p\}.
\]
For any $r\in\{1,\dots,n\}$, the geodesic $C_{r,1}$ has its endpoints at $e^{\sfrac{2i\pi(r-1)}{n}}$ and $e^{\Psfrac{2i\pi (r-1)}{n}+\Psfrac{2i\pi}{np}}$, and the other $C_{r,s}$ are obtained by rotating $C_{r,1}$ successively by angle $\sfrac{2\pi}{np}$ about the origin. The geodesics $C_{r,s}$ induce a partition of the arc $J_r$ into $p$ arcs $J_{r,1},\dots, J_{r,p}$, where $J_{r,s}$ is the arc of $\mathbb{S}^1$ of length $\sfrac{2\pi}{np}$ connecting the endpoints of $C_{r,s}$.

The bi-infinite geodesics $C_{r,s}$, $r\in\{1,\dots,n\},\ s\in\{1,\dots,p\}$, bound a closed ideal $np$-gon (in the topology of $\D$), which we call $\pmb{\Pi}$. We~will now introduce Möbius maps of the disk that pair the sides of $\pmb{\Pi}$. To do so, we~will exploit the symmetry $M_\omega$ of~$\pmb{\Pi}$. Specifically, we~will prescribe the side-pairings for $C_{1,1},\dots, C_{1,p}$ explicitly, and conjugate these side-pairing transformations by powers of $M_\omega$ to define pairings for the other sides of $\pmb{\Pi}$. Let us denote the diameter of $\mathbb{S}^1$ with endpoints at $\pm e^{\sfrac{i\pi}{n}}$ by $\ell$. Now observe that the Möbius map $g_{1,s}$ obtained by post-composing the reflection in $C_{1,s}$ with the reflection in $\ell$ carries $C_{1,s}$ to $C_{1,p+1-s}$. In~particular, $g_{1,p+1-s}=g_{1,s}^{-1}$. Note that when $p$ is odd, then $g_{1,\psfrac{p+1}{2}}$ is an involution with a fixed point on $C_{1,\psfrac{p+1}{2}}$.

By the Poincaré polygon theorem, the Möbius maps
\[
g_{r,s}:=M_\omega^{r-1}\circ g_{1,s}\circ M_{\omega}^{-(r-1)},\quad r\in\{1,\dots,n\},\ s\in\{1,\dots,p\}
\]
generate a Fuchsian group $\pmb{\Gamma_{n,p}}$, and $\pmb{\Pi}$ is a closed fundamental domain for the $\pmb{\Gamma_{n,p}}$\nobreakdash-action on $\D$. Moreover, $\faktor{\D}{\pmb{\Gamma_{n,p}}}$ is biholomorphic to

\begin{itemize}
\item
a sphere with $\Psfrac{np}{2}+1$ punctures for $p$ even, and

\item
a sphere with $\Psfrac{n(p-1)}{2}+1$ punctures and $n$ order two orbifold points for $p$ odd.
\end{itemize}

\skpt
\begin{remark}\label{reconcile_rem_1}
\begin{enumerate}
\item
The integer $p$ can be thought of as the number of `pockets' in each sector of angular width $2\pi/n$. On the other hand, the integer $n$ plays the role of $\nu$ appearing in the definition of the class $\cS$ of genus zero orbifolds (see Section~\ref{intro_sec}).

\item
When $n\geq 3$, the orbifold $\faktor{\D}{\pmb{\Gamma_{n,p}}}$ is an $n$-fold cyclic cover of a base genus zero orbifold $\pmb{\Sigma}\in\cS$ with $\lfloor p/2\rfloor +1$ punctures, zero/one order two orbifold point depending on the parity of $p$, and an order $n$ orbifold point.
\end{enumerate}
\end{remark}

We now look at the Bowen-Series map $A_{\pmb{\Gamma_{n,p}}}^{\mathrm{BS}}$ equipped with the above fundamental domain $\pmb{\Pi}$ and side-pairing transformations (\cf \cite{bowen_series}). By~definition, the map
\[
A_{\pmb{\Gamma_{n,p}}}^{\mathrm{BS}}:\overline{\D}\setminus\Int{\pmb{\Pi}}\to \overline{\D}
\]
acts as $g_{r,s}$ on the closure of the hyperbolic half-plane enclosed by the geodesic $C_{r,s}$ and the arc $J_{r,s}$.
It is now easily checked that $A_{\pmb{\Gamma_{n,p}}}^{\mathrm{BS}}$ is continuous on $\mathbb{S}^1\setminus\sqrt[n]{1}$, and the left and right-hand limits of $A_{\pmb{\Gamma_{n,p}}}^{\mathrm{BS}}$ at the points of $\sqrt[n]{1}$ lie in the set $\sqrt[n]{1}$. We~will now use this fact to pass to a factor of $A_{\pmb{\Gamma_{n,p}}}^{\mathrm{BS}}$ that is continuous everywhere.

Consider the bordered (orbifold) Riemann surfaces
\[
\mathcal{Q}:=\faktor{\overline{\D}}{\langle M_\omega\rangle},\quad \mathcal{Q}_1:=\faktor{\left(\overline{\D}\setminus \Int{\pmb{\Pi}}\right)}{\langle M_\omega\rangle},
\]
and note that a closed fundamental domain for the action of $\langle M_\omega \rangle$ on $\overline{\D}$ is given by
\[
\{z\in\overline{\D}:\ 0\leq\arg{z}\leq\sfrac{2\pi}{n}\}\cup\{0\}.
\]
Thus, $\mathcal{Q}$ is biholomorphic to the surface obtained from the above fundamental domain by identifying the radial line segments $\{r:0\leq r\leq 1\}$ and $\{re^{\sfrac{2\pi i}{n}}:0\leq r\leq 1\}$ by~$M_\omega$.

By construction, the Bowen-Series map $A_{\pmb{\Gamma_{n,p}}}^{\mathrm{BS}}$ commutes with $M_\omega$, and hence it can be pushed forward via the quotient map $\mathfrak{q}:\overline{\D}\to\mathcal{Q}$ to define a map
\[
\mathfrak{q}\circ A_{\pmb{\Gamma_{n,p}}}^{\mathrm{BS}}\circ\mathfrak{q}^{-1}: \mathcal{Q}_1\to\mathcal{Q}.
\]
Note that the map $z\mto z^{n}$ yields a conformal isomorphism $\xi$ between the (bordered) surfaces $\mathcal{Q}$ and $\overline{\D}$. Finally, we~set
\[
A_{\pmb{\Gamma_{n,p}}}^{\mathrm{fBS}}:= \xi\circ \bigl(\mathfrak{q}\circ A_{\pmb{\Gamma_{n,p}}}^{\mathrm{BS}}\circ\mathfrak{q}^{-1}\bigr) \circ \xi^{-1}: \mathcal{D}_{\pmb{\Gamma_{n,p}}}:=\xi(\mathcal{Q}_1)\to \overline{\D}.
\]
Note that $A_{\pmb{\Gamma_{n,p}}}^{\mathrm{fBS}}:\mathbb{S}^1\to\mathbb{S}^1$ is an orientation-preserving covering map of degree $np-1$. By~\cite[Lem.\,3.7]{LMMN}, the map $A_{\pmb{\Gamma_{n,p}}}^{\mathrm{fBS}}\vert_{\mathbb{S}^1}$ is expansive. Moreover, $A_{\pmb{\Gamma_{n,p}}}^{\mathrm{fBS}}$ has $p$ critical points at $\xi(\mathfrak{q}(g_{1,s}(0)))$, $s\in\{1,\dots, p\}$, and each of them has multiplicity $n-1$. All these critical points are mapped to $0$.

\begin{definition}\label{factor_bs_base_def}
We call the map $A_{\pmb{\Gamma_{n,p}}}^{\mathrm{fBS}}: \mathcal{D}_{\pmb{\Gamma_{n,p}}}\to \overline{\D}$ the \emph{factor Bowen-Series} map of $\pmb{\Gamma_{n,p}}$ equipped with the fundamental domain $\pmb{\Pi}$.
\end{definition}

\begin{remark}\label{rmk-Pi}
It is worth noting that
$A_{\pmb{\Gamma_{n,p}}}^{\mathrm{fBS}}$ extends continuously to the boundary $\partial \, \pmb{\Pi}/\langle M_\omega\rangle$ of the ideal triangle $\pmb{\Pi}/\langle M_\omega\rangle$. Further, this action on $\partial \, \pmb{\Pi}/\langle M_\omega\rangle$ is by complex conjugation. In~what follows, it will be important that the action of $A_{\pmb{\Gamma_{n,p}}}^{\mathrm{fBS}}$ on $\partial \, \pmb{\Pi}/\langle M_\omega\rangle$ is by an involution.
\end{remark}

\subsection{Deformations and moduli spaces of factor Bowen-Series maps}\label{factor_bs_gen_subsec}

The preferred fundamental domain $\pmb{\Pi}$ of $\pmb{\Gamma_{n,p}}$ can be used to equip every marked Fuchsian group in the Teichmüller space $\mathrm{Teich}(\pmb{\Gamma_{n,p}})$ of $\pmb{\Gamma_{n,p}}$ with a preferred fundamental domain. This, in turn, will allow us to define factor Bowen-Series maps for all marked groups in a suitable symmetry locus of $\mathrm{Teich}(\pmb{\Gamma_{n,p}})$.

Recall that the Teichmüller space $\mathrm{Teich}(\pmb{\Gamma_{n,p}})$ is the space of Möbius conjugacy classes of discrete, faithful, strongly type-preserving representations of $\pi_1\left(\faktor{\D}{\pmb{\Gamma_{n,p}}}\right)\cong \pmb{\Gamma_{n,p}}$ into $\mathrm{Aut}(\D)\cong \mathrm{PSL}_2(\R)$. Any such representation $\rho: \pmb{\Gamma_{n,p}}\to\Gamma$ is given by $\rho(g)=\psi_\rho\circ g\circ\psi_\rho^{-1},\ g\in\pmb{\Gamma_{n,p}}$, where $\psi_\rho$ is a quasiconformal homeomorphism of $\widehat{\C}$ that preserves $\D$.
The quasiconformal homeomorphism is not unique. However,
two such quasiconformal homeomorphisms inducing the same representation $\rho$ agree on $\mathbb{S}^1$. We~can and will choose a quasiconformal homeomorphism $\psi_\rho$ such that $\psi_\rho(1)=1$, and $\psi_\rho(\pmb{\Pi})$ is a hyperbolic ideal polygon.
The ideal polygon $\psi_\rho(\pmb{\Pi})$ is a preferred fundamental domain for the marked group $\Gamma$.

We denote by $\mathrm{Teich}^\omega(\pmb{\Gamma_{n,p}})$ the collection of $(\rho: \pmb{\Gamma_{n,p}}\to\Gamma)\in\mathrm{Teich}(\pmb{\Gamma_{n,p}})$ that commute with conjugation by $M_\omega$; \ie
\[
\rho(M_\omega\circ g\circ M_\omega^{-1})=M_\omega\circ \rho(g)\circ M_\omega^{-1},\quad g\in\pmb{\Gamma_{n,p}}.
\]
This is equivalent to requiring that the associated quasiconformal map $\psi_\rho$ commutes with $M_\omega$.

For each $(\rho: \pmb{\Gamma_{n,p}}\to\Gamma)\in \mathrm{Teich}^\omega(\pmb{\Gamma_{n,p}})$, consider the Bowen-Series map
\[
A_{\Gamma}^{\mathrm{BS}} \equiv\psi_\rho\circ A_{\pmb{\Gamma_{n,p}}}^{\mathrm{BS}}\circ\psi_\rho^{-1}:\overline{\D}\setminus \Int{\psi_\rho(\pmb{\Pi})}\to\overline{\D}
\]
associated with the marked group $\Gamma$ equipped with the preferred fundamental domain $\psi_\rho(\pmb{\Pi})$. By~definition, $A_\Gamma^{\mathrm{BS}}$ commutes with $M_\omega$, and thus can be pushed forward via the quotient map $\mathfrak{q}:\overline{\D}\to\mathcal{Q}$. As~in the previous section, this gives rise to a map
\[
A_{\Gamma}^{\mathrm{fBS}}:= \xi\circ \left(\mathfrak{q}\circ A_{\Gamma}^{\mathrm{BS}} \circ\mathfrak{q}^{-1}\right) \circ \xi^{-1}: \mathcal{D}_{\Gamma}:= \xi(\mathfrak{q}(\overline{\D}\setminus \Int{\psi_\rho(\pmb{\Pi})}))\to \overline{\D},
\]
that is quasiconformally conjugate to $A_{\pmb{\Gamma_{n,p}}}^{\mathrm{fBS}}$.

\begin{definition}\label{factor_bs_gen_def}
For $(\rho: \pmb{\Gamma_{n,p}}\to\Gamma)\in \mathrm{Teich}^\omega(\pmb{\Gamma_{n,p}})$ induced by the quasiconformal map $\psi_\rho$, the map $A_{\Gamma}^{\mathrm{fBS}}: \mathcal{D}_{\Gamma}\to \overline{\D}$ is called the \emph{factor Bowen-Series} map of $\Gamma$ equipped with the fundamental domain $\psi_\rho(\pmb{\Pi})$.
\end{definition}

\begin{remark}
For the symmetric base group $\pmb{\Gamma_{n,p}}$, the generators chosen in Section~\ref{factor_bs_base_subsec} are compositions of two anti-Möbius reflections. However, the (quasiconformally deformed) preferred generators for $(\rho: \pmb{\Gamma_{n,p}}\to\Gamma)\in \mathrm{Teich}^\omega(\pmb{\Gamma_{n,p}})$ in general do not admit such symmetries.
\end{remark}

The group $\widehat{\pmb{\Gamma}}_{n,p}$ generated by $\pmb{\Gamma_{n,p}}$ and $M_\omega$ is an index $n$ extension of $\pmb{\Gamma_{n,p}}$. Clearly, the set
\[
\widehat{\pmb{\Pi}}:=\{ z\in\pmb{\Pi}:\ 0\leq\arg{z}\leq\sfrac{2\pi}{n}\}\cup\{0\}
\]
is a closed fundamental domain for the action of $\widehat{\pmb{\Gamma}}_{n,p}$ on $\D$. It~follows that $\faktor{\D}{\pmb{\widehat{\Gamma}}_{n,p}}$ is biholomorphic to
\begin{itemize}
\item
a sphere with $\sfrac{p}{2}+1$ punctures and an order $n$ orbifold point for $p$ even, and

\item
a sphere with $\psfrac{p+1}{2}$ punctures, an order two orbifold point and an order $n$ orbifold point for $p$ odd.
\end{itemize}

The space of factor Bowen-Series maps constructed above is parametrized by $\mathrm{Teich}^\omega(\pmb{\Gamma_{n,p}})$, which in turn can be identified with the Teichmüller space $\mathrm{Teich}(\widehat{\pmb{\Gamma}}_{n,p})$. Specifically, for the representation $(\rho: \pmb{\Gamma_{n,p}}\to\Gamma)\in \mathrm{Teich}^\omega(\pmb{\Gamma_{n,p}})$, we~define
\[
\widehat{\Gamma}:=\langle\Gamma, M_\omega\rangle,
\]
and associate with $\rho$ the representation
\[
(\widehat{\rho}: \widehat{\pmb{\Gamma}}_{n,p}\to\widehat{\Gamma})\in \mathrm{Teich}(\widehat{\pmb{\Gamma}}_{n,p}),\quad \text{where}\quad \widehat{\rho}\vert_{\pmb{\Gamma_{n,p}}}\equiv\rho\vert_{\pmb{\Gamma_{n,p}}}\quad \text{and}\quad \widehat{\rho}(M_\omega)=M_\omega.
\]
Thus, $\mathrm{Teich}^\omega(\pmb{\Gamma_{n,p}})$ is the same as the Teichmüller space of the orbifold group
$\widehat{\pmb{\Gamma}}_{n,p}=\pmb{\Gamma_{n,p}} \rtimes \langle M_\omega\rangle$.

\skpt
\begin{remark}\label{reconcile_rem_2}
\begin{enumerate}
\item
The orbifold $\faktor{\D}{\widehat{\pmb{\Gamma}}_{n,p}}$ is a base genus zero orbifold $\pmb{\Sigma}\in\cS$ with $\lfloor p/2\rfloor +1$ punctures, zero/one order two orbifold point depending on the parity of $p$, and an order $n$ orbifold point when $n\geq 3$ (\cf Remark~\ref{reconcile_rem_1}). Thus, any $\Sigma\in\mathrm{Teich}(\pmb{\Sigma})$ is uniformized by some Fuchsian group $\Gamma\in\mathrm{Teich}(\widehat{\pmb{\Gamma}}_{n,p})$.
\medskip

\item
The fact that the chosen fundamental domain and side-pairings of the base group $\pmb{\Gamma_{n,p}}$ admit a $2\pi/n$ rotation symmetry and that the representations \hbox{$(\rho: \pmb{\Gamma_{n,p}}\to\Gamma)\in\mathrm{Teich}^\omega(\pmb{\Gamma_{n,p}})$} respect this symmetry, together imply that the orbifolds $\faktor{\D}{\Gamma}$ have an order $n$ isometry. Quotienting by this isometry yields an $n$-fold (branched) covering $\faktor{\D}{\Gamma}\to\faktor{\D}{\widehat{\Gamma}}$.
\end{enumerate}
\end{remark}

\begin{figure}[htb]
\vspace*{-2\baselineskip}
\[
\begin{tikzcd}
\faktor{\D}{\Gamma} \arrow{d}[swap]{\mathrm{cover}} \arrow[r,rightsquigarrow] & A_{\Gamma}^{\mathrm{BS}} \arrow{d}{\mathrm{factor}} \\
\faktor{\D}{\widehat{\Gamma}} & A_{\Gamma}^{\mathrm{fBS}}
\end{tikzcd}
\]
\vspace*{-\baselineskip}
\caption{Covering orbifolds and factor Bowen-Series maps}
\label{going_up_down_2_fig}
\end{figure}

We summarize the main properties of factor Bowen-Series maps below.

\skpt
\begin{proposition}\label{factor_bs_prop}
\begin{enumerate}
\item\label{factor_bs_prop1} $A_{\Gamma}^{\mathrm{fBS}}:\mathbb{S}^1\to\mathbb{S}^1$ is a piecewise analytic, orientation-preserving, expansive, covering map of degree $np-1$. In~particular, it is topologically conjugate to $z^{np-1}\vert_{\mathbb{S}^1}$; \ie there exists a homeomorphism $\mathfrak{h}_\Gamma: {\mathbb{S}^1} \to {\mathbb{S}^1}$ such that $\mathfrak{h}_\Gamma (z^{np-1}) = \big( A_{\Gamma}^{\mathrm{fBS}}(\mathfrak{h}_\Gamma (z))\big)$.

\item\label{factor_bs_prop2} The restriction $A_{\Gamma}^{\mathrm{fBS}}: \left(A_{\Gamma}^{\mathrm{fBS}}\right)^{-1}(\mathcal{D}_{\Gamma})\to \mathcal{D}_{\Gamma}$ is a covering map of degree $np-1$. In~particular, it maps each component of $\left(A_{\Gamma}^{\mathrm{fBS}}\right)^{-1}(\Int{\mathcal{D}_{\Gamma}})$ homeomorphically onto some component of $\Int{\mathcal{D}_{\Gamma}}$.

\item\label{factor_bs_prop3} The restriction $A_{\Gamma}^{\mathrm{fBS}}: \left(A_{\Gamma}^{\mathrm{fBS}}\right)^{-1}(\overline{\D}\setminus\mathcal{D}_{\Gamma})\to \overline{\D}\setminus\mathcal{D}_{\Gamma}$ has degree $np$. If $n\geq 2$, there are $p$ critical points of $A_{\Gamma}^{\mathrm{fBS}}$, each of multiplicity $n-1$, in $\left(A_{\Gamma}^{\mathrm{fBS}}\right)^{-1}(\overline{\D}\setminus\mathcal{D}_{\Gamma})$.
All these critical points are mapped to the unique critical value $0$ of $A_{\Gamma}^{\mathrm{fBS}}$.

\item\label{factor_bs_prop4} The set $\overline{\D}\setminus\mathcal{D}_\Gamma$ is (the interior of) a topological $p$-gon with its $p$ vertices on~$\mathbb{S}^1$. Each such vertex is a fixed point or a $2$-periodic point under $A_{\Gamma}^{\mathrm{fBS}}$.

\item\label{factor_bs_prop5} For $\Gamma=\pmb{\Gamma_{n,p}}$, the factor Bowen-Series map acts as complex conjugation on $\partial\mathcal{D}_{\pmb{\Gamma_{n,p}}}\cap\D$, which is the boundary of the ideal polygon $\overline{\D}\setminus\mathcal{D}_{\pmb{\Gamma_{n,p}}}$. For a general $\Gamma$, the factor Bowen-Series map $A_{\Gamma}^{\mathrm{fBS}}$ acts by an involution on $\partial\mathcal{D}_{\Gamma}\cap\D$.
\end{enumerate}
\end{proposition}

\begin{proof}
It is enough to verify the assertions for the base map $A_{\pmb{\Gamma_{n,p}}}^{\mathrm{fBS}}$.
We endow the bordered Riemann surface $\mathcal{\mathcal{Q}}$ with a preferred choice of complex coordinates via its identification with $\{z\in\overline{\D}:\ 0\leq\arg{z}\leq\sfrac{2\pi}{n}\}\cup\{0\}$ (with the boundary radial lines glued together).

\eqref{factor_bs_prop1} The facts that $A_{\pmb{\Gamma_{n,p}}}^{\mathrm{BS}}$ is continuous on $\mathbb{S}^1\setminus\sqrt[n]{1}$ and the left and right-hand limits of $A_{\pmb{\Gamma_{n,p}}}^{\mathrm{BS}}$ at the points of $\sqrt[n]{1}$ lie in the set $\sqrt[n]{1}$ together imply that
\[
\overline{A}_{\pmb{\Gamma_{n,p}}}^{\mathrm{BS}}:=\mathfrak{q}\circ A_{\pmb{\Gamma_{n,p}}}^{\mathrm{BS}}\circ\mathfrak{q}^{-1}
\]
is continuous.

Let us partition each arc $J_{1,s}\subset\mathbb{S}^1$, $s\in\{1,\dots,p\}$, into sub-arcs $J_{1,s}^1,\dots, J_{1,s}^{m(s)}$, where each $J_{1,s}^i$ is a connected component of some $g_{1,s}^{-1}(J_r)$, with $i\in\{1,\dots,m(s)\}$, $r\in\nobreak\{1,\dots,n\}$. Then, with the above choice of coordinates on $\mathcal{Q}$,
\begin{itemize}
\item $\overline{A}_{\pmb{\Gamma_{n,p}}}^{\mathrm{BS}}$ acts as a Möbius map $h_{i,s}$ (called a \emph{piece} of $\overline{A}_{\pmb{\Gamma_{n,p}}}^{\mathrm{BS}}$) on $\mathfrak{q}(J_{1,s}^i)$; specifically, $h_{i,s}$ is a composition of $g_{1,s}$ with a power of~$M_\omega$,

\item $h_{i,s}(\mathfrak{q}(J_{1,s}^i))$ is the union of finitely many sub-arcs from the collection
\[
\{\mathfrak{q}(J_{1,s}):s\in\{1,\dots,p\}\}.
\]
\end{itemize}
It follows that with the above choice of coordinates on $\mathcal{Q}$, the map $\overline{A}_{\pmb{\Gamma_{n,p}}}^{\mathrm{BS}}$ is a piecewise Möbius, orientation-preserving, covering map of $\partial\mathcal{Q}$. The statement that $\overline{A}_{\pmb{\Gamma_{n,p}}}^{\mathrm{BS}}:\partial\mathcal{Q}\to\partial\mathcal{Q}$ has degree $np-1$ follows from the fact that all but finitely many points in $\mathbb{S}^1$ have $np-1$ preimages under the map $A_{\pmb{\Gamma_{n,p}}}^{\mathrm{BS}}$ (since $A_{\pmb{\Gamma_{n,p}}}^{\mathrm{BS}}$ maps each arc $J_{r,s}$ to $\mathbb{S}^1\setminus\nobreak\Int{J_{r,p+1-s}}$). Finally, expansivity of $\overline{A}_{\pmb{\Gamma_{n,p}}}^{\mathrm{BS}}\vert_{\partial\mathcal{Q}}$ is a consequence of the fact that each $g_{1,s}$ has derivative larger than one on $\Int{J_{1,s}}$, for $s\in\{1,\dots,p\}$ (\cf \cite[Lem.\,3.8]{LMMN}).

Since $\xi:\mathcal{Q}\to\overline{\D}$ is a biholomorphism, the properties of $\overline{A}_{\pmb{\Gamma_{n,p}}}^{\mathrm{BS}}$ listed in the previous paragraph imply that the map $A_{\pmb{\Gamma_{n,p}}}^{\mathrm{fBS}}\equiv \xi\circ \overline{A}_{\pmb{\Gamma_{n,p}}}^{\mathrm{BS}}\vert_{\partial\mathcal{Q}}\circ\xi^{-1}:\mathbb{S}^1\to\mathbb{S}^1$ is a piecewise analytic (\emph{not} piecewise Möbius when $n>1$), orientation-preserving, expansive, covering map of degree $np-1$.

The existence of the circle homeomorphism $\mathfrak{h}_\Gamma$ that conjugates $z^{np-1}$ to $A_\Gamma^{\mathrm{fBS}}$ now follows from the fact that two expansive circle coverings of the same degree are topologically conjugate (\cf \cite[Property~(2'), p.\,99]{CR80}).

\eqref{factor_bs_prop2} and \eqref{factor_bs_prop3} The Bowen-Series map $A_{\pmb{\Gamma_{n,p}}}^{\mathrm{BS}}$ sends the hyperbolic half-plane bounded by the geodesic $C_{1,s}$ and the arc $J_{1,s}\subset\mathbb{S}^1$ conformally to the complement of the hyperbolic half-plane bounded by the geodesic $C_{1,p+1-s}$ and the arc $J_{1,p+1-s}\subset\mathbb{S}^1$, $s\in\{1,\dots,p\}$. Hence, the region $\overline{\D}\setminus\mathcal{D}_{\Gamma}$ is covered $np$ times by $A_{\pmb{\Gamma_{n,p}}}^{\mathrm{fBS}}$, while $\mathcal{D}_{\Gamma}$ is covered $np-1$ times. Clearly, $A_{\Gamma}^{\mathrm{fBS}}: \left(A_{\Gamma}^{\mathrm{fBS}}\right)^{-1}(\mathcal{D}_{\Gamma})\to \mathcal{D}_{\Gamma}$ is piecewise conformal, and hence a covering map.

To locate the critical points of $A_{\pmb{\Gamma_{n,p}}}^{\mathrm{fBS}}$, let us denote the union of the radial lines in~$\overline{\D}$ at angles $\sfrac{2j\pi}{n}$, $j\in\{0,\dots,n-1\}$, by $\mathcal{P}$. Note that $A_{\pmb{\Gamma_{n,p}}}^{\mathrm{fBS}}$ maps each $\xi(\mathfrak{q}(g_{1,s}^{-1}(\mathcal{P})))$ to the line segment $[0,1]$ and sends $\xi(\mathfrak{q}(g_{1,s}^{-1}(0)))$ to $0$, for $s\in\{1,\dots,p\}$. It~follows that for each $s\in\{1,\dots,p\}$, the point $\xi(\mathfrak{q}(g_{1,s}^{-1}(0)))$ is a critical point of multiplicity $n-1$ with associated critical value $0$.

\eqref{factor_bs_prop4} and \eqref{factor_bs_prop5} These follow from Definitions~\ref{factor_bs_base_def} and~\ref{factor_bs_gen_def}, and the discussion in Section~\ref{factor_bs_base_subsec}.
\end{proof}

We denote the circle homeomorphism that conjugates $z^{np-1}$ to $A_{\pmb{\Gamma_{n,p}}}^{\mathrm{fBS}}$ by $\mathfrak{h}_0$, normalized such that $\mathfrak{h}_0$ sends the fixed point $1$ of $z^{np-1}$ to the fixed point $1$ of $A_{\pmb{\Gamma_{n,p}}}^{\mathrm{fBS}}$. Further, the quasiconformal conjugacy $\psi_\rho$ between $\pmb{\Gamma_{n,p}}$ and $\Gamma$ induces a quasiconformal conjugacy $\widehat{\psi}_\rho$ between $A_{\pmb{\Gamma_{n,p}}}^{\mathrm{fBS}}$ and $A_{\Gamma}^{\mathrm{fBS}}$.
In the remainder of the paper, we~will work with the normalized conjugacy $\mathfrak{h}_\Gamma=\widehat{\psi}_\rho\circ\mathfrak{h}_0:\mathbb{S}^1\to\mathbb{S}^1$ between $z^{np-1}\vert_{\mathbb{S}^1}$ and~$A_{\Gamma}^{\mathrm{fBS}}\vert_{\mathbb{S}^1}$.

Let us denote the set of $p$ ideal boundary points of $\overline{\D}\setminus \mathcal{D}_{\Gamma}$ on $\mathbb{S}^1$ by $S_{\Gamma}$.
Under the circle homeomorphism $\mathfrak{h}_\Gamma$, the set $S_{\Gamma}$ is pulled back to the set
\[
\displaystyle\mathcal{A}_p:=\{\sfrac{i}{p}:i\in\{0,\dots,p-1\}\}
\]
(here we identify $\mathbb{S}^1$ with $\R/\Z$). See Proposition~\ref{factor_bs_prop} (parts 4, 5). More precisely, if~$p$ is even (respectively, odd), the two points (respectively, one point) of $S_{\Gamma}$ that are (respectively, is) fixed by $A_{\Gamma}^{\mathrm{fBS}}$ correspond to $0,\sfrac{p}{2p}=\sfrac{1}{2}$ (respectively, corresponds to~$0$), and the $2$-cycles of $A_{\Gamma}^{\mathrm{fBS}}$ in $S_{\Gamma}$ correspond to the $2$-cycles $\pm \sfrac{i}{p}$, $i\in\{1,\dots,\lfloor \frac{p-1}{2}\rfloor\}$, of $m_{np-1}$.


\section{Conformal matings of factor Bowen-Series maps with polynomials}\label{conf_mating_sec}

The goal of this section is to prove Theorem~\href{https://jep.centre-mersenne.org/articles/10.5802/jep.315/}{A}. In~fact, we~will prove a more general statement that allows the polynomials to lie in arbitrary hyperbolic components in the connectedness locus.

\subsection{The notion of conformal mating}\label{conf_mating_def_subsec}

Let $n, p$ be two positive integers with \hbox{$np\geq 3$}, and $P$ be a monic, centered complex polynomial of degree $d:=np-1$ with a connected and locally connected Julia set $\mathcal{J}(P)$. Let $\mathcal{K}(P)$ denote the filled Julia set.
Recall that there exists a unique conformal map
\[
\psi_P:\widehat{\C}\setminus\overline{\D}\to\mathcal{B}_\infty(P):=\widehat{\C}\setminus\mathcal{K}(P)
\]
(called the \emph{Böttcher coordinate} of $P$) that conjugates $z^{d}$ to $P$, and is tangent to the identity map near infinity (\cf \cite[Th.\,9.1]{Mil06}). As~$\mathcal{J}(P)$ is locally connected, $\psi_P$~extends continuously to $\mathbb{S}^1$ to yield a semi-conjugacy between $z^d\vert_{\mathbb{S}^1}$ and $P\vert_{\mathcal{J}(P)}$.

We now define the notion of topological/conformal mating of $P$ and $A_{\Gamma}^{\mathrm{fBS}}$, where $\Gamma\in\mathrm{Teich}^\omega(\pmb{\Gamma_{n,p}})$.
Recall from Section~\ref{factor_bs_gen_subsec} that there exists a normalized homeomorphism $\mathfrak{h}_\Gamma:\mathbb{S}^1\to\mathbb{S}^1$ that conjugates $z^d$ to $A_{\Gamma}^{\mathrm{fBS}}$.
Let us now consider the disjoint union $\mathcal{K}(P)\sqcup \overline{\D}$ and the map\vspace*{-3pt}
\begin{gather*}
P\sqcup A_{\Gamma}^{\mathrm{fBS}}: \mathcal{K}(P)\sqcup \mathcal{D}_{\Gamma}\to \mathcal{K}(P)\sqcup \overline{\D},\\
\bigl(P\sqcup A_{\Gamma}^{\mathrm{fBS}}\bigr)\vert_{\mathcal{K}(P)}=P,\quad \bigl(P\sqcup A_{\Gamma}^{\mathrm{fBS}}\bigr)\vert_{\mathcal{D}_{\Gamma}}=A_{\Gamma}^{\mathrm{fBS}}.
\end{gather*}
Let $\sim_{\mathrm{m}}$ (here `m' stands for mating) be the equivalence relation on $\mathcal{K}(P)\sqcup \overline{\D}$ generated~by\vspace*{-3pt}
\begin{equation}
\psi_{P}(z)\sim_{\mathrm{m}} \mathfrak{h}_\Gamma(\overline{z})),\quad \text{for all}\ z\in\mathbb{S}^1.
\label{conf_mat_equiv_rel}
\end{equation}
The map $P\sqcup A_{\Gamma}^{\mathrm{fBS}}$ descends to a continuous map $P\mate A_{\Gamma}^{\mathrm{fBS}}$ to the quotient\vspace*{-3pt}
\[
\mathcal{K}(P)\ \mate\ \overline{\D}\ :=\ \left(\mathcal{K}(P)\sqcup \overline{\D}\right)/\sim_{\mathrm{m}}\ \cong\ \mathbb{S}^2.
\]
The map $P\mate A_{\Gamma}^{\mathrm{fBS}}$ is called the \emph{topological mating} of $P$ and $A_{\Gamma}^{\mathrm{fBS}}$.
We say that $P$ and $A_{\Gamma}^{\mathrm{fBS}}$ are \emph{conformally mateable} if the topological $2$-sphere $\mathcal{K}(P)\mate \overline{\D}$ admits a complex structure that turns the topological mating $P\mate A_{\Gamma}^{\mathrm{fBS}}$ into a holomorphic~map.

Here is an equivalent formulation (\cf \cite[Def.\,4.14]{petersen-meyer-mating}).
\begin{definition}\label{conf_mat_def}
The maps $P$ and $A_{\Gamma}^{\mathrm{fBS}}$ are conformally mateable if there exist a continuous map $F\colon \mathrm{Dom}(F)\subsetneq\widehat{\C}\to\widehat{\C}$ (called a \emph{conformal mating} of $A_{\Gamma}^{\mathrm{fBS}}$ and $P$) that is complex-analytic in the interior of $\mathrm{Dom}(F)$ and continuous maps\vspace*{-3pt}
\[
\mathfrak{X}_P:\mathcal{K}(P)\to\widehat{\C}\quad \text{and}\quad \mathfrak{X}_\Gamma: \overline{\D}\to\widehat{\C},
\]
conformal on $\Int{\mathcal{K}(P)}$ and $\D$ (respectively), satisfying
\begin{enumerate}
\item\label{topo_cond} $\mathfrak{X}_P\left(\mathcal{K}(P)\right)\cup \mathfrak{X}_\Gamma\left(\overline{\D}\right) = \widehat{\C}$,

\item\label{dom_cond} $\mathrm{Dom}(F)= \mathfrak{X}_P(\mathcal{K}(P))\cup\mathfrak{X}_\Gamma(\mathcal{D}_{\Gamma})$,

\item $\mathfrak{X}_P\circ P(z) = F\circ \mathfrak{X}_P(z),\quad \mathrm{for}\ z\in\mathcal{K}(P)$,

\item $\mathfrak{X}_\Gamma\circ A_{\Gamma}^{\mathrm{fBS}}(w) = F\circ \mathfrak{X}_\Gamma(w),\quad \mathrm{for}\ w\in
\mathcal{D}_{\Gamma}$,\quad and

\item\label{identifications} $\mathfrak{X}_P(z)=\mathfrak{X}_\Gamma(w)$ if and only if $z\sim_{\mathrm{m}} w$ where $\sim_{\mathrm{m}}$ is the equivalence relation on $\mathcal{K}(P)\sqcup \overline{\D}$ defined by Relation~\eqref{conf_mat_equiv_rel}.
\end{enumerate}
The semi-conjugacies $\mathfrak{X}_P, \mathfrak{X}_\Gamma$ are called the \emph{mating semi-conjugacies} associated with the conformal mating $F$ of $P$ and $A_{\Gamma}^{\mathrm{fBS}}$. When mating semi-conjugacies are injective, they are simply referred to as \emph{mating conjugacies}.
\end{definition}

\subsection{Existence of conformal matings}\label{conf_mat_exists_subsec}

Let $(\rho: \pmb{\Gamma_{n,p}}\to\Gamma)\in\mathrm{Teich}^\omega(\pmb{\Gamma_{n,p}})$ and~$P$ a~monic, centered, hyperbolic complex polynomial of degree $d=np-1$ with a connected Julia set. We~now state and prove a generalization of \cite[Th.\,3.6]{MM1}.


\begin{theorem}\label{conf_mat_thm}
There exists a conformal mating $F$ of $P:\mathcal{K}(P)\to\mathcal{K}(P)$ and $A_{\Gamma}^{\mathrm{fBS}}:\mathcal{D}_{\Gamma}\to \overline{\D}$. Moreover, $F$ is unique up to Möbius conjugacy.
\end{theorem}

We will start with a technical lemma.

\begin{definition}\label{david_def}
An orientation-preserving homeomorphism $H: U\to V$ between domains in the Riemann sphere $\widehat{\C}$ is called a \textit{David homeomorphism} if it lies in the Sobolev class $W^{1,1}_{\text{loc}}(U)$ and there exist constants $C,\alpha,\varepsilon_0>0$ with
\begin{align}\label{david_cond}
\sigma(\{z\in U: |\mu_H(z)|\geq 1-\varepsilon\}) \leq Ce^{-\alpha/\varepsilon}, \quad \varepsilon\leq \varepsilon_0.
\end{align}
\end{definition}
Here $\sigma$ is the spherical measure, and
$\mu_H= \frac{\partial H/ \partial\overline{z}}{\partial H/\partial z}$
is the Beltrami coefficient of $H$ (see \cite[Chap.\,20]{AIM09} for more background on David homeomorphisms).

\begin{lemma}\label{david_ext_lem}
The circle homeomorphism $\mathfrak{h}_\Gamma$ of Proposition~\ref{factor_bs_prop} continuously extends to a David homeomorphism of $\D$.
\end{lemma}
\begin{proof}
Recall from Section~\ref{factor_bs_gen_subsec} that $h_\Gamma=\widehat{\psi}_\rho\circ\mathfrak{h}_0$, where $\widehat{\psi}_\rho$ is the restriction of a global quasiconformal map and $\mathfrak{h}_0$ conjugates $z^{np-1}$ to $A_{\pmb{\Gamma_{n,p}}}^{\mathrm{fBS}}$. Hence, it suffices to check that $\mathfrak{h}_0$ continuously extends to a David homeomorphism of $\D$.

We will use the notation introduced in Proposition~\ref{factor_bs_prop}. In~particular, we~endow~$\mathcal{\mathcal{Q}}$ with a preferred choice of complex coordinates via its identification with the set \hbox{$\{z\in\overline{\D}:\ 0\leq\arg{z}\leq\sfrac{2\pi}{n}\}\cup\{0\}$} (with the boundary radial lines glued together).

The partition of $\partial\mathcal{Q}$ into the arcs
\[
\bigl\{\mathfrak{q}(J_{1,s}^i): \, s\in\{1,\dots,p\},\ i\in\{1,\dots,m(s)\}\bigr\}
\]
does \emph{not} necessarily give a Markov partition for $\overline{A}_{\pmb{\Gamma_{n,p}}}^{\mathrm{BS}}$ since the map may send both endpoints of such a partition piece to $1$. However, we~can refine the above partition by pulling it back under $\overline{A}_{\pmb{\Gamma_{n,p}}}^{\mathrm{BS}}$, and this produces a Markov partition~$\{I_k\}$.

Each piece $h_{i,s}\vert_{I_k}$ of $\overline{A}_{\pmb{\Gamma_{n,p}}}^{\mathrm{BS}}$ extends conformally as $h_{i,s}$ to a neighborhood of $I_k$ in~$\widetilde{\mathcal{Q}}$, where $\widetilde{\mathcal{Q}}$ is the double of $\mathcal{Q}$ and $I_k\subset \mathfrak{q}(J_{1,s}^i)$.
Finally, since the pieces of $\overline{A}_{\pmb{\Gamma_{n,p}}}^{\mathrm{BS}}$ are Möbius (with respect to the preferred coordinates on $\mathcal{Q}$), which send round disks to round disks, we~can choose round disk neighborhoods $U_k\subset \widetilde{\mathcal{Q}}$ of the interiors of the Markov partition pieces $\Int{I_k}$ (intersecting $\partial\mathcal{Q}$ orthogonally) such that if \hbox{$\overline{A}_{\pmb{\Gamma_{n,p}}}^{\mathrm{BS}}(I_k)\supset I_{k'}$}, then $\overline{A}_{\pmb{\Gamma_{n,p}}}^{\mathrm{BS}}(U_k)\supset U_{k'}$.

The properties of $\overline{A}_{\pmb{\Gamma_{n,p}}}^{\mathrm{BS}}$ listed in the previous paragraph imply that the map $A_{\pmb{\Gamma_{n,p}}}^{\mathrm{fBS}}\equiv \xi\circ \overline{A}_{\pmb{\Gamma_{n,p}}}^{\mathrm{BS}}\vert_{\partial\mathcal{Q}}\circ\xi^{-1}:\mathbb{S}^1\to\mathbb{S}^1$ is a piecewise analytic orientation-preserving expansive covering map of degree $d\geq2$ admitting a Markov partition $\{\xi(I_k)\}$ satisfying conditions (4.1) and (4.2) of \cite[Th.\,4.12]{LMMN}. Moreover, each periodic breakpoint of its piecewise analytic definition is symmetrically parabolic (\cf \cite[Def.\,4.6, Rem.\,4.7]{LMMN}). By~\cite[Th.\,4.12]{LMMN}, the orientation-preserving homeomorphism $\mathfrak{h}_0:\mathbb{S}^1\to\mathbb{S}^1$ that conjugates the map $z^d\vert_{\mathbb{S}^1}$ to $A_{\pmb{\Gamma_{n,p}}}^{\mathrm{fBS}}\vert_{\mathbb{S}^1}$ (and sends $1$ to $1$) extends continuously as a David homeomorphism of~$\D$.
\end{proof}

\begin{proof}[Proof of Theorem~\ref{conf_mat_thm}]
As $\mathcal{J}(P)$ is locally connected, $\psi_P$ extends to a continuous surjection $\psi_P:\mathbb{S}^1\to\mathcal{J}(P)$ semi-conjugating $z^d$ to $P$. Also note that since $P$ is hyperbolic with connected Julia set, $\mathcal{B}_\infty(P)$ is a John domain and $\mathcal{J}(P)$ is removable for $W^{1,1}$ functions \cite[Th.\,4]{JS00}.

Let $\mathfrak{h}_\Gamma:\overline{\D}\to\overline{\D}$ be a continuous extension of $\mathfrak{h}_\Gamma:\mathbb{S}^1\to\mathbb{S}^1$ that is a David homeomorphism on $\D$. The existence of such an extension is guaranteed by Lemma~\ref{david_ext_lem}.
Also recall that $\mathfrak{h}_\Gamma$ conjugates $z^d\vert_{\mathbb{S}^1}$ to $A_{\Gamma}^{\mathrm{fBS}}\vert_{\mathbb{S}^1}$.
Consider the topological dynamical system
\begin{equation*}
\widetilde{F}(w):=\begin{cases}
P & \text{on}\ \mathcal{K}(P), \\
\psi_P\circ\eta\circ\mathfrak{h}_\Gamma^{-1}\circ A_{\Gamma}^{\mathrm{fBS}}\circ\mathfrak{h}_\Gamma\circ\eta\circ\psi_P^{-1} & \text{on}\ \psi_P\left(\eta\left(\mathfrak{h}_\Gamma^{-1}\left(\mathcal{D}_{\Gamma}\right)\right)\right)\subset\mathcal{B}_\infty(P),
\end{cases}
\end{equation*}
where $\eta(z)=1/z$. By~equivariance properties of $\mathfrak{h}_\Gamma:\mathbb{S}^1\to\mathbb{S}^1$ and $\psi_P:\mathbb{S}^1\to\mathcal{J}(P)$, the two definitions agree on $\mathcal{J}(P)$. We~denote the domain of $\widetilde{F}$ by $\mathrm{Dom}(\widetilde{F})$.

We define a Beltrami coefficient $\mu$ on the sphere as follows. On $\mathcal{K}(P)$ we set $\mu$ to be
the standard complex structure. On $\mathcal{B}_\infty(P)$, we~set $\mu$ to be the pullback of the standard complex structure (on $\D$) under the map $\mathfrak{h}_\Gamma\circ\eta\circ\psi_P^{-1}$. Since $\mathfrak{h}_\Gamma\circ\eta\circ\psi_P^{-1}$ is a David homeomorphism (by~\cite[Prop.\,2.5(iv)]{LMMN}), it follows that $\mu$ is a David coefficient on $\widehat{\C}$. It~is easy to check that $\mu$ is $\widetilde{F}$-invariant.

The David integrability theorem (see \cite{David}, \cite[Th.\,20.6.2, p.\,578]{AIM09}) provides us with a David homeomorphism $\mathfrak{H}:\widehat{\C}\to\widehat{\C}$ such that the pullback of the standard complex structure under $\mathfrak{H}$ is equal to $\mu$. Conjugating $\widetilde{F}$ by $\mathfrak{H}$, we~obtain the map
\[
F:=\mathfrak{H}\circ\widetilde{F}\circ\mathfrak{H}^{-1}:\mathfrak{H}(\mathrm{Dom}(\widetilde{F}))\to\widehat{\C}.
\]
We set $\mathrm{Dom}(F):= \mathfrak{H}(\mathrm{Dom}(\widetilde{F}))$.

We proceed to show that $F$ is holomorphic on $\Int{\mathrm{Dom}(F)}$.
As $\mathcal{J}(P)$ is removable for functions in class $W^{1,1}$, it follows from \cite[Th.\,2.7]{LMMN} that $\mathfrak{H}(\mathcal{J}(P))$ is locally conformally removable. Hence, it suffices to show that $F$ is holomorphic on the interior of $\mathrm{Dom}(F)\setminus\mathfrak{H}(\mathcal{J}(P))$. Indeed, this would imply that the continuous map $F$ is holomorphic on $\Int{\mathrm{Dom}(F)}$ away from the finitely many critical points of $F$. One can then conclude that $F$ is holomorphic on $\Int{\mathrm{Dom}(F)}$ using the Riemann removability theorem.

To this end, first observe that both the maps $\mathfrak{h}_\Gamma\circ\eta\circ\psi_P^{-1}$ and $\mathfrak{H}$ are David homeomorphisms on $\mathcal{B}_\infty(P)$ straightening $\mu\vert_{\mathcal{B}_\infty(P)}$. By~\cite[Th.\,20.4.19, p.\,565]{AIM09}, the map $\mathfrak{h}_\Gamma\circ\eta\circ\psi_P^{-1}\circ\mathfrak{H}^{-1}$ is conformal on $\mathfrak{H}(\mathcal{B}_\infty(P))$. It~now follows from the definitions of $\widetilde{F}$ and $F$ that $F$ is holomorphic on $\mathfrak{H}(\mathcal{B}_\infty(P))\cap\Int{\mathrm{Dom}(F)}$. Similarly, both the identity map and the map $\mathfrak{H}$ are David homeomorphisms on each component of $\Int{\mathcal{K}(P)}$ straightening $\mu$. Once again by \cite[Th.\,20.4.19, p.\,565]{AIM09}, $\mathfrak{H}$ is conformal on each component of $\Int{\mathcal{K}(P)}$. By~definition of $\widetilde{F}$ and $F$, it now follows that $F$ is holomorphic on each interior component of $\mathfrak{H}(\mathcal{K}(P))$.
This completes the proof of the fact that $F$ is holomorphic on the interior of $\mathrm{Dom}(F)$.

Finally, we~set $\mathfrak{X}_P:=\mathfrak{H}:\mathcal{K}(P)\to\widehat{\C}$ and $\mathfrak{X}_\Gamma:=\mathfrak{H}\circ\psi_P\circ\eta\circ\mathfrak{h}_\Gamma^{-1}:\overline{\D}\to\widehat{\C}$. It~is readily checked that these maps satisfy the requirements of Definition~\ref{conf_mat_def}. Thus, $F$ is a conformal mating of $P$ and $A_{\Gamma}^{\mathrm{fBS}}$.

Now suppose that there is another conformal mating $F_1$ of $P$ and $A_{\Gamma}^{\mathrm{fBS}}$. Then the respective mating semi-conjugacies paste together to yield a homeomorphism of $\widehat{\C}$ which is conformal away from $\mathfrak{H}(\mathcal{J}(P))$ and conjugates $F$ to $F_1$. Conformal removability of $\mathfrak{H}(\mathcal{J}(P))$ now implies that this homeomorphism is a Möbius map; \ie $F$ and $F_1$ are Möbius conjugate.
\end{proof}

\begin{thebibliography}{99}
\bibitem[MM23a]{MM1} M.Mj and S.Mukherjee, ``Combining rational maps and Kleinian groups via orbit equivalence'', \emph{Proc. London Math. Soc.}(3)\textbf{126}(2023),no.5,1740--1809.
\bibitem[BS79]{bowen_series} R.Bowen and C.Series, ``Markov maps associated with Fuchsian groups'', \emph{Publ. Math. Inst. Hautes \'Etudes Sci.}(1979),no.50,153--170.
\bibitem[LMMN25]{LMMN} M.Lyubich,S.Merenkov,S.Mukherjee and D.Ntalampekos, \emph{David extension of circle homeomorphisms,welding,mating,and removability}, Mem. Amer. Math. Soc.313,no.1588, American Mathematical Society,Providence,RI,2025.
\bibitem[CR80]{CR80} E.M.Coven and W.L.Reddy, ``Positively expansive maps of compact manifolds'', in \emph{Global theory of dynamical systems}(Evanston,Ill.,1979), Lect. Notes in Math.819, Springer,Berlin,1980,96--110.
\bibitem[Mil06]{Mil06} J.Milnor, \emph{Dynamics in one complex variable}, third ed., Annals of Math. Studies160, Princeton University Press,Princeton,NJ,2006.
\bibitem[PM12]{petersen-meyer-mating} C.L.Petersen and D.Meyer, ``On the notions of mating'', \emph{Ann. Fac. Sci. Toulouse Math.}(6)\textbf{21}(2012),no.5,839--876.
\bibitem[AIM09]{AIM09} K.Astala,T.Iwaniec and G.Martin, \emph{Elliptic partial differential equations and quasiconformal mappings in the plane}, Princeton Math. Series48, Princeton University Press,Princeton,NJ,2009.
\bibitem[JS00]{JS00} P.W.Jones and S.K.Smirnov, ``Removability theorems for Sobolev functions and quasiconformal maps'', \emph{Ark. Mat.}\textbf{38}(2000),no.2,263--279.
\bibitem[Dav88]{David} G.David, ``Solutions de l'équation de Beltrami avec $\|\mu\|_\infty=1$'', \emph{Ann. Acad. Sci. Fenn. Ser.AI Math.}\textbf{13}(1988),no.1,25--70.
\end{thebibliography}
\end{document}
